%------------ Debut Packages ----------------
\usepackage[T1]{fontenc}
\usepackage[utf8]{inputenc}
\usepackage{lmodern}
\usepackage[french]{babel}
\usepackage{xcolor} % pour ajouter de la couleur (si besoin)
\usepackage{multirow}
\usepackage{makecell}
\usepackage{booktabs}
\usepackage{float}
\usepackage{rotating}
\usepackage{tablefootnote}
\usepackage{hyperref}
\usepackage{bookmark}
\usepackage{pdfpages}
\usepackage{lipsum}
\usepackage{listings}
\usepackage{subcaption}
\usepackage{url}

\lstdefinelanguage{Shell}{
  keywords={sudo, docker, nmap, npm, node},
  sensitive=true,
  comment=[l]{\#},
  morestring=[b]",
  morestring=[b]'
}

\lstset{
  language=Shell,
  basicstyle=\ttfamily\small,
  numbers=none,
  backgroundcolor=\color{gray!10},
  showspaces=false,
  showstringspaces=false,
  frame=single,
  breaklines=true,
  captionpos=b,
  keywordstyle=\color{blue}\bfseries,
  stringstyle=\color{teal},
  commentstyle=\color{gray}
}

% Pour personnaliser les entêtes et pieds de pages
\usepackage{fancyhdr}

\pagestyle{fancy}
\renewcommand{\chaptermark}[1]{\markboth{\chaptername \ \thechapter.\ #1}{}}
\renewcommand{\sectionmark}[1]{\markright{\thesection.\ #1}}
\renewcommand{\thesection}{\arabic{section}}

%------------ Fin Packages ----------------
